\section{Theoretical Background}
\label{sec:theoretical-background}

This chapter defines the physical configuration, governing equations, and
turbulence statistics required for the present work. It first establishes the
incompressible, non-buoyant channel-flow baseline; thermal stratification will
be introduced later.

\subsection{Governing Equations and Physical Configuration}
\label{sec:channel-flow-governing-equations}

Plane channel flow consists of fluid motion between two stationary, parallel
walls separated by the full channel height $2\delta$, where $\delta$ is the
channel half-height. The coordinates $x$, $y$, and $z$ denote the streamwise,
spanwise, and wall-normal directions, respectively. The corresponding
instantaneous velocity components are $u$, $v$, and $w$, and the velocity
vector is $\mathbf{u}=(u,v,w)$. Figure~\ref{fig:channel-geometry} shows the
coordinate convention used in this chapter.

\begin{figure}[htbp]
  \centering
  \setlength{\unitlength}{1mm}
  \begin{picture}(120,46)
    \put(12,8){\line(1,0){96}}
    \put(12,36){\line(1,0){96}}
    \put(15,3){stationary wall, $z=-\delta$}
    \put(15,39){stationary wall, $z=\delta$}
    \put(43,22){\vector(1,0){30}}
    \put(75,20){$x$}
    \put(43,22){\vector(0,1){11}}
    \put(39,33){$z$}
    \put(43,22){\circle*{1.2}}
    \put(38,17){$y$ out of the page}
    \put(86,22){\vector(1,0){15}}
    \put(88,25){mean flow}
  \end{picture}
  \caption{Plane-channel configuration: $x$ is streamwise, $y$ is spanwise
  (out of the page), and $z$ is wall-normal.}
  \label{fig:channel-geometry}
\end{figure}

\subsubsection{Conservation of Mass}

For an incompressible fluid of constant density, conservation of mass reduces
to
\begin{equation}
  \boldsymbol{\nabla}\mathbin{\cdot}\mathbf{u}=0,
  \label{eq:incompressible-continuity}
\end{equation}
where $\boldsymbol{\nabla}$ is the spatial gradient operator. Equation
\eqref{eq:incompressible-continuity} requires the velocity field to be
divergence-free and expresses local conservation of fluid mass.

\subsubsection{Conservation of Momentum}

The present formulation treats the fluid as Newtonian, with constant density
$\rho$ and constant kinematic viscosity $\nu$. Conservation of momentum is
then described by the incompressible Navier--Stokes equations,
\begin{equation}
  \frac{\partial \mathbf{u}}{\partial t}
  +(\mathbf{u}\mathbin{\cdot}\boldsymbol{\nabla})\mathbf{u}
  =-\frac{1}{\rho}\boldsymbol{\nabla}p
  +\nu\boldsymbol{\nabla}^{2}\mathbf{u},
  \label{eq:incompressible-navier-stokes}
\end{equation}
where $t$ is time, $p$ is pressure, and $\boldsymbol{\nabla}^{2}$ is the
Laplace operator. In SI units, $\rho$, $\nu$, and $p$ are measured in
$\mathrm{kg\,m^{-3}}$, $\mathrm{m^2\,s^{-1}}$, and pascals, respectively.
The two terms on the left-hand side represent local and convective
acceleration; the right-hand side contains pressure acceleration and viscous
momentum diffusion \cite{pope2000}.

\subsubsection{Pressure-Gradient Forcing}

A streamwise pressure gradient maintains the channel flow. The pressure is
written as a linear mean contribution plus a periodic remainder
$\widetilde{p}$,
\begin{equation}
  p(x,y,z,t)=-Gx+\widetilde{p}(x,y,z,t),
  \qquad
  G=-\frac{\mathrm{d}\langle p\rangle}{\mathrm{d}x}>0,
  \label{eq:channel-pressure-decomposition}
\end{equation}
where $G$ is the imposed streamwise pressure drop per unit length in
$\mathrm{Pa\,m^{-1}}$, and angle brackets denote a statistical average.
Substitution into Eq.~\eqref{eq:incompressible-navier-stokes} gives
\begin{equation}
  \frac{\partial \mathbf{u}}{\partial t}
  +(\mathbf{u}\mathbin{\cdot}\boldsymbol{\nabla})\mathbf{u}
  =-\frac{1}{\rho}\boldsymbol{\nabla}\widetilde{p}
  +\nu\boldsymbol{\nabla}^{2}\mathbf{u}
  +\frac{G}{\rho}\mathbf{e}_x,
  \label{eq:pressure-driven-navier-stokes}
\end{equation}
where $\mathbf{e}_x$ is the unit vector in the positive streamwise direction.
The final term is the uniform acceleration associated with the mean pressure
gradient. Equation~\eqref{eq:pressure-driven-navier-stokes} defines the
physical forcing considered here; its implementation by the solver is not
specified at this stage.

\subsubsection{Computational Domain and Boundary Conditions}

The walls are located at $z=-\delta$ and $z=\delta$. Impermeability and the
no-slip condition require
\begin{equation}
  \mathbf{u}(x,y,-\delta,t)
  =\mathbf{u}(x,y,\delta,t)=\mathbf{0}.
  \label{eq:channel-no-slip}
\end{equation}
The zero vector in Eq.~\eqref{eq:channel-no-slip} states that all three
velocity components vanish at each stationary wall. For a channel domain of
streamwise length $L_x$ and spanwise length $L_y$, periodicity in the two
wall-parallel directions is expressed as
\begin{equation}
  \mathbf{u}(x+L_x,y,z,t)=\mathbf{u}(x,y,z,t),
  \qquad
  \mathbf{u}(x,y+L_y,z,t)=\mathbf{u}(x,y,z,t).
  \label{eq:channel-periodicity}
\end{equation}
These conditions represent a repeating section of the channel and make $x$
and $y$ statistically homogeneous. The instantaneous flow remains
three-dimensional and time-dependent \cite{kimMoinMoser1987}.

\subsubsection{Fully Developed Channel Flow}

Fully developed channel flow is defined statistically: its mean quantities do
not vary in the homogeneous directions. The mean velocity therefore has the
form
\begin{equation}
  \langle\mathbf{u}\rangle=\bigl(U(z),0,0\bigr),
  \label{eq:fully-developed-channel-flow}
\end{equation}
where $U(z)=\langle u\rangle$ is the mean streamwise velocity. Equation
\eqref{eq:fully-developed-channel-flow} describes the mean field only and does
not imply that the instantaneous turbulent field is steady.

\subsection{Turbulence Statistics for Validation}
\label{sec:reynolds-decomposition-turbulence-statistics}

Validation of a turbulent channel calculation is based on converged
statistics rather than individual instantaneous fields. The quantities needed
for the present work are defined below using the coordinate convention of
Subsection~\ref{sec:channel-flow-governing-equations}.

\subsubsection{Reynolds Decomposition and Averaging}

The instantaneous velocity is separated into mean and fluctuating parts by the
Reynolds decomposition,
\begin{equation}
  u_i(\mathbf{x},t)=U_i(\mathbf{x})+u_i'(\mathbf{x},t),
  \qquad
  U_i=\langle u_i\rangle,
  \qquad i\in\{1,2,3\},
  \label{eq:reynolds-decomposition}
\end{equation}
where $\mathbf{x}=(x,y,z)$ is the position vector, $u_i$ is an instantaneous
velocity component, $U_i$ is its mean, $u_i'$ is its fluctuation, and $i$ is
the component index. The assignments $u_1=u$, $u_2=v$, and $u_3=w$ correspond
to the streamwise, spanwise, and wall-normal components. In component form,
$u=U+u'$, $v=V+v'$, and $w=W+w'$, where $U=\langle u\rangle$,
$V=\langle v\rangle$, and $W=\langle w\rangle$. All velocity quantities have
units of $\mathrm{m\,s^{-1}}$. For the fully developed flow in
Eq.~\eqref{eq:fully-developed-channel-flow}, $U=U(z)$ and $V=W=0$. The
fluctuations satisfy
\begin{equation}
  \langle u_i'\rangle=0,
  \qquad
  \langle u'\rangle=\langle v'\rangle=\langle w'\rangle=0.
  \label{eq:zero-mean-velocity-fluctuations}
\end{equation}
Equation~\eqref{eq:zero-mean-velocity-fluctuations} follows directly from the
definition of the mean \cite{pope2000}.

For an instantaneous flow quantity $\phi$, the ideal time average is
\begin{equation}
  \langle\phi\rangle_t(x,y,z)
  =\lim_{T\rightarrow\infty}\frac{1}{T}
  \int_{t_0}^{t_0+T}\phi(x,y,z,t)\,\mathrm{d}t,
  \label{eq:turbulence-time-average}
\end{equation}
where $t_0$ is the start of the sampling interval and $T$ is its duration in
seconds. In a statistically stationary calculation, the converged result is
independent of $t_0$. Averaging additionally over the homogeneous directions
gives
\begin{equation}
  \langle\phi\rangle(z)
  =\frac{1}{L_xL_y}\int_0^{L_y}\int_0^{L_x}
  \langle\phi\rangle_t(x,y,z)\,\mathrm{d}x\,\mathrm{d}y.
  \label{eq:channel-space-time-average}
\end{equation}
Equation~\eqref{eq:channel-space-time-average} is the statistical average used
below. For finite simulation data, the sampling interval must be long enough
for each reported quantity to converge. Under an ergodic interpretation, this
space--time average estimates the corresponding ensemble average. Statistical
stationarity refers to the statistics, not to a steady instantaneous field.

The second moments are collected in the kinematic Reynolds-stress tensor,
\begin{equation}
  R_{ij}(z)=\langle u_i'u_j'\rangle,
  \label{eq:reynolds-stress-tensor}
\end{equation}
where $R_{ij}$ has units of $\mathrm{m^2\,s^{-2}}$. The diagonal components
$R_{11}$, $R_{22}$, and $R_{33}$ are the streamwise, spanwise, and wall-normal
normal stresses; the off-diagonal components describe correlations between
different velocity fluctuations.

\subsubsection{RMS Velocity Fluctuations}

The root-mean-square (RMS) velocity fluctuations are
\begin{equation}
  u_{\mathrm{rms}}(z)=\sqrt{\langle u'^2\rangle},
  \qquad
  v_{\mathrm{rms}}(z)=\sqrt{\langle v'^2\rangle},
  \qquad
  w_{\mathrm{rms}}(z)=\sqrt{\langle w'^2\rangle}.
  \label{eq:rms-velocity-fluctuations}
\end{equation}
These quantities are the square roots of $R_{11}$, $R_{22}$, and $R_{33}$,
respectively, and have units of $\mathrm{m\,s^{-1}}$. Their profiles measure
the typical streamwise, spanwise, and wall-normal fluctuation magnitudes across
the channel and are standard quantities for comparison with reference DNS data
\cite{kimMoinMoser1987}.

\subsubsection{Reynolds Shear Stress}

The Reynolds-stress component relevant to wall-normal transport of streamwise
momentum is
\begin{equation}
  -R_{13}(z)=-\langle u'w'\rangle.
  \label{eq:reynolds-shear-stress}
\end{equation}
Here, $u'$ and $w'$ are the streamwise and wall-normal velocity fluctuations.
The kinematic stress in Eq.~\eqref{eq:reynolds-shear-stress} has units of
$\mathrm{m^2\,s^{-2}}$; multiplication by $\rho$ gives the dimensional
turbulent shear stress $-\rho\langle u'w'\rangle$ in pascals. This correlation
quantifies turbulent momentum transfer between wall-parallel layers and is a
central validation quantity for channel flow \cite{pope2000}.

\subsubsection{Turbulent Kinetic Energy}

The turbulent kinetic energy per unit mass is
\begin{equation}
  k(z)=\frac{1}{2}\langle u_i'u_i'\rangle
  =\frac{1}{2}\left(
    \langle u'^2\rangle+\langle v'^2\rangle+\langle w'^2\rangle
  \right),
  \label{eq:turbulent-kinetic-energy}
\end{equation}
where a repeated index denotes summation over the three Cartesian components.
The quantity $k$ is the mean kinetic energy of the velocity fluctuations per
unit fluid mass. Its SI units are $\mathrm{m^2\,s^{-2}}$, equivalently
$\mathrm{J\,kg^{-1}}$.

\subsubsection{Validation Quantities for the Present Work}

The present work will later extract the mean streamwise velocity $U(z)$, the
three RMS velocity fluctuations, the Reynolds shear stress
$-\langle u'w'\rangle$, and the turbulent kinetic energy $k(z)$ from BRINE.
Their definitions are independent of the solver implementation, but the
coordinate mapping and sampling procedure must be confirmed before the output
is processed. Agreement with trusted non-buoyant channel DNS data is required
before differences caused by buoyancy are interpreted.

\subsection{Wall Scaling, Friction Reynolds Number, and Law of the Wall}
\label{sec:wall-scaling-law-wall}

Wall scaling relates the near-wall velocity field to the shear exerted by a
solid boundary. For the mean streamwise velocity $U(z)$ defined in
Eq.~\eqref{eq:fully-developed-channel-flow}, the positive magnitude of the wall
shear stress is
\begin{equation}
  \tau_w
  =\mu\left|\frac{\mathrm{d}U}{\mathrm{d}z}\right|_w
  =\rho\nu\left|\frac{\mathrm{d}U}{\mathrm{d}z}\right|_w,
  \label{eq:wall-shear-stress}
\end{equation}
where the subscript $w$ denotes evaluation at either wall, $\mu=\rho\nu$ is
the dynamic viscosity, and $\nu$ is the kinematic viscosity introduced in
Eq.~\eqref{eq:incompressible-navier-stokes}. The absolute value removes the
opposite gradient signs at the lower and upper walls. Physically, $\tau_w$ is
the streamwise force transmitted per unit wall area; its SI unit is the pascal,
$\mathrm{Pa}=\mathrm{N\,m^{-2}}$.

The corresponding friction velocity, channel half-height, and friction
Reynolds number are
\begin{equation}
  u_\tau=\sqrt{\frac{\tau_w}{\rho}},
  \qquad
  h=\frac{L_z}{2},
  \qquad
  Re_\tau=\frac{u_\tau h}{\nu}.
  \label{eq:friction-scaling-quantities}
\end{equation}
Here, $u_\tau$ has units of velocity but is a shear-based scaling quantity, not
the local fluid velocity. The half-height $h$ is identical to $\delta$ in the
centred geometry of Fig.~\ref{fig:channel-geometry}, and $L_z$ is the full
wall-normal channel length. The parameter $Re_\tau$ measures the separation
between the outer length $h$ and the near-wall viscous length; increasing
$Re_\tau$ generally increases the range of dynamically relevant wall-turbulence
scales \cite{pope2000}.

The planned project convention uses $L_z=2$, $|\Pi|=1$, and $\rho=1$, where
$\Pi=\mathrm{d}\langle p\rangle/\mathrm{d}x=-G$ is the signed mean pressure
gradient associated with Eq.~\eqref{eq:channel-pressure-decomposition}. For a
fully developed plane channel, the mean streamwise momentum balance gives
$\tau_w=h|\Pi|$. Consequently,
\begin{equation}
  h=1,
  \qquad
  \tau_w=1,
  \qquad
  u_\tau=1,
  \qquad
  \nu=\mu=\frac{1}{Re_\tau}
  \quad (\rho=1).
  \label{eq:project-wall-scaling-normalization}
\end{equation}
Equation~\eqref{eq:project-wall-scaling-normalization} documents the intended
nondimensional setup. It is not a result from a completed simulation.

Let $d_w$ be the non-negative distance from a point to the nearest wall. In the
centred coordinate system of Fig.~\ref{fig:channel-geometry},
$d_w=h-|z|$; for data stored on $0\leq z\leq L_z$, the equivalent expression is
$d_w=\min(z,L_z-z)$. The viscous length, wall distance, and mean velocity in
wall units are
\begin{equation}
  \ell_\nu=\frac{\nu}{u_\tau},
  \qquad
  z^+=\frac{d_w}{\ell_\nu}=\frac{d_wu_\tau}{\nu},
  \qquad
  U^+=\frac{U}{u_\tau}.
  \label{eq:viscous-length-wall-units}
\end{equation}
The RMS quantities from Eq.~\eqref{eq:rms-velocity-fluctuations} and the
Reynolds shear stress from Eq.~\eqref{eq:reynolds-shear-stress} are scaled as
\begin{equation}
  \begin{aligned}
    u_{\mathrm{rms}}^+&=\frac{u_{\mathrm{rms}}}{u_\tau},
    &v_{\mathrm{rms}}^+&=\frac{v_{\mathrm{rms}}}{u_\tau},
    &w_{\mathrm{rms}}^+&=\frac{w_{\mathrm{rms}}}{u_\tau},\\
    -\langle u'w'\rangle^+
    &=-\frac{\langle u'w'\rangle}{u_\tau^2}
    =\frac{-\rho\langle u'w'\rangle}{\tau_w}.
  \end{aligned}
  \label{eq:inner-scaled-turbulence-statistics}
\end{equation}

BRINE uses $z$ and $w$ for the wall-normal coordinate and velocity component,
respectively. Much of the wall-turbulence literature instead uses $y$ as the
wall-normal coordinate and therefore denotes wall distance by $y^+$. That
conventional symbol may be retained in comparison figures only when it is
defined explicitly as distance from the nearest wall in wall units. It must not
be confused with the BRINE coordinate $y$, which is spanwise, or with the
spanwise velocity $v$. The same mapping converts a literature stress
$-\langle u'v'\rangle$ into the BRINE quantity $-\langle u'w'\rangle$ when the
reference uses $v$ for wall-normal velocity.

The near-wall flow is commonly described using four approximate regions. In
the viscous sublayer, molecular viscosity dominates the mean momentum transfer.
The buffer layer is a transition in which viscous and turbulent effects are
both important. Farther from the wall, the logarithmic or overlap region links
inner and outer scaling, whereas the outer region is influenced by the channel
half-height, pressure-gradient forcing, and centreline conditions. These
regions do not have universal sharp boundaries; their extent depends in part
on Reynolds number \cite{pope2000}.

For a smooth wall, the leading inner relation in the viscous sublayer is
\begin{equation}
  U^+=z^+.
  \label{eq:viscous-law-of-wall}
\end{equation}
In the logarithmic overlap region, the generic law of the wall is
\begin{equation}
  U^+=\frac{1}{\kappa_{\mathrm{vK}}}\ln(z^+)+B,
  \label{eq:logarithmic-law-of-wall}
\end{equation}
where $\kappa_{\mathrm{vK}}$ is the dimensionless von K\'arm\'an constant and
$B$ is the additive intercept. The symbol $\kappa_{\mathrm{vK}}$ is distinct
from the thermal diffusivity $\kappa$, which has units of
$\mathrm{m^2\,s^{-1}}$. No project-specific numerical values are assigned to
$\kappa_{\mathrm{vK}}$ or $B$ at this stage \cite{pope2000}.

Wall-unit profiles provide a common basis for comparing channel calculations
at a specified $Re_\tau$. The planned unstratified benchmark at $Re_\tau=590$
will therefore compare $U^+$, the three RMS fluctuations in
Eq.~\eqref{eq:inner-scaled-turbulence-statistics}, and the inner-scaled
Reynolds shear stress with established channel-DNS statistics
\cite{kimMoinMoser1987,moserKimMansour1999}. This comparison extends the
validation quantities
identified in Subsection~\ref{sec:reynolds-decomposition-turbulence-statistics};
it defines a future validation procedure and does not imply that validation has
already been completed.
